\documentclass[10pt,letterpaper]{article}
\usepackage[utf8]{inputenc}
\usepackage[T1]{fontenc}
\usepackage[margin=0.65in]{geometry}
\usepackage{amsmath,amssymb}
\usepackage{booktabs}
\usepackage{longtable}
\usepackage{graphicx}
\usepackage{hyperref}
\hypersetup{colorlinks=true,linkcolor=blue,citecolor=blue,urlcolor=blue}

\title{\textbf{Mosquito Tree of Life: Triassic--Cenozoic Culicidae Phylogenomics (242.0 to 0.0 Ma; Crown Culicidae 145.0 Ma), Single Origin of Hematophagy, Salivary D7 \& Apyrase Recruitment, and Secondary Loss in Toxorhynchites}}
\author{\textbf{Bryan P. White} \\ CRV Science Labs \\ \texttt{bryan@crvscience.com}}
\date{November 7, 2026 \quad $\cdot$ \quad Preprint ID: PREPRINT-14 (v3.0.0 Empirical Preprint) \quad $\cdot$ \quad DOI: \href{https://doi.org/10.5281/zenodo.24039858}{10.5281/zenodo.24039858}}

\begin{document}
\maketitle

\begin{abstract}
\noindent Resolving the Mosquito Tree of Life (Diptera: Culicidae) and testing whether vertebrate hematophagy evolved once along the family stem or independently across subfamilies confront severe third-codon mitochondrial AT saturation (mean 78.84\% AT overall and $>91\%$ AT at third-codon positions across Culicidae mitogenomes) versus balanced nuclear coding composition (44.33\%--51.06\% AT). Here we execute a dual-pipeline empirical phylogenomic synthesis across 120 primary verified NCBI RefSeq and GenBank sequence records totaling 1,648,476 bp (1,642,055 bp primary FASTA on disk across 100 complete mitochondrial genomes [\texttt{T001}--\texttt{T100} and Internal Nodes \texttt{N01}--\texttt{N98} = 1,595,551 bp] plus 20 core RefSeq nuclear transcripts [\texttt{NUC-101}--\texttt{NUC-120} = 46,504 bp raw on disk; 52,925 bp across 48 verified RefSeq orthologs in 6 nuclear families]), aligning 100\% of all downloaded sequences across their full lengths in both \textbf{MAFFT v7.526} (\textbf{30,983 aligned columns} across 100 mitogenomes in 61.27 s + \textbf{7,390 aligned columns} across 20 nuclear loci in 5.42 s = \textbf{38,373 bp} full-length Total-Evidence supermatrix in 66.69 s, alongside the 16,032 bp and 4,140 bp codon-partitioned supermatrices) and Native Rust \textbf{\texttt{gotoh-msa}} (\textbf{89,830 aligned columns} in 0.96 s; $64.1\times$ faster than MAFFT), evaluated across a 32-chain multi-model combinatorial optimization sweep using \textbf{IQ-TREE v2.4.0} ($B = 1,000$ UFBoot2 replicates), \textbf{IQ-TREE LSD2} (\texttt{--date}) relaxed clock dating, and \textbf{PAML v4.10.10} (\texttt{codeml}). Across the full-length 30,983 bp MAFFT alignment, Hybrid \texttt{RY3} \texttt{GTR+F+I+G4} ($\ln L = -436,799.68$, $\text{BIC} = 875,739.99$, 187.82 s) outperforms External Raw \texttt{GTR+F+I+G4} ($\ln L = -495,670.85$) by $+58,871.17$ nats, while across the full-length 89,830 bp Native Rust \texttt{gotoh-msa} alignment, Native \texttt{RY3} \texttt{GTR+F+R2} ($\ln L = -615,839.29$, $\text{BIC} = 1,234,016.74$, 132.88 s) outperforms Native Raw \texttt{GTR+F+R2} ($\ln L = -694,172.90$) by $+78,333.61$ nats. On the codon-partitioned matrices, \texttt{RY3} reduces third-codon RCFV by 76.3\% ($0.0737 \rightarrow 0.0175$), achieving $\ln L = -152,134.21$ ($\text{BIC} = 306,476.00$) on the 16,032 bp total-evidence supermatrix (Pipeline A) and $\ln L = -67,577.15$ ($\text{BIC} = 136,961.58$; mean UFBoot2 $= 90.4\% \pm 1.6\%$; 100\% UFBoot2 on all major subfamily/tribe stems) on the 4,140 bp 100-species mitochondrial supermatrix (Pipeline B). LSD2 relaxed-clock dating places the Culicomorpha root in the Middle Triassic (242.00 Ma), the Chaoboridae + Culicidae sister split in the Early Jurassic (198.72 Ma), and Crown Culicidae at the Jurassic--Cretaceous boundary (145.00 Ma; 95\% CI: 145.00--160.23 Ma; 100\% UFBoot2), followed by Crown Anophelinae (117.59 Ma) and Crown Culicinae (125.00 Ma). Verified obligate non-biting \textit{Toxorhynchites} (\texttt{T060}, \texttt{T062}, \texttt{T063}) nests inside Culicinae (stem 109.45 Ma; crown 53.36 Ma; 100\% UFBoot2), while PAML v4.10.10 3-branch modeling on authentic \textit{Toxorhynchites rutilus} RefSeq mRNAs (\texttt{GCF\_029784165.1}) indicates $2.11\times$ relaxed purifying selection on salivary \textit{D7L1} ($\omega = 0.2138$ vs.\ $0.1011$) and $2.09\times$--$4.48\times$ relaxed purifying selection on secreted salivary apyrase ($\omega = 0.1025$ vs.\ $0.0229$--$0.0491$) alongside $1.57\times$ tighter purifying selection on \textit{Vitellogenin-A1} ($\omega = 0.0929$ vs.\ $0.1457$, $p < 10^{-16}$). Finally, our Dual-Pass Blinded Validation ($Z_{\text{blind}} = +16.41$ SD vs.\ $Z_{\text{ctx}} = +16.10$ SD; $D_{\text{ctx}} = 0.31$ SD) uncovers a published GenBank sample-swap error in \texttt{PQ585799.1} (deposited as \textit{Toxorhynchites rutilus}, but 100.0000\% identical across 14,500 bp to \textit{Culex tarsalis} \texttt{PQ585801.1}).
\end{abstract}

\section{Biological Context \& Dual-Pipeline Empirical Architecture}
This study reconstructs the \textbf{Mosquito Tree of Life} (Diptera: Culicidae) under a Full Step-Up Trajectory (Tier 1 Exploratory $\rightarrow$ Tier 2 Validation $\rightarrow$ Tier 3 Scientific) in Dual-Pass Blinded Mode~\cite{White2026RuNIE}. Across 120 primary verified NCBI RefSeq and GenBank sequence records totaling 1,648,476 bp (100 complete Culicidae and Culicomorpha mitochondrial genomes across Tip Nodes \texttt{T001}--\texttt{T100} and Internal Nodes \texttt{N01}--\texttt{N98} = 1,595,551 bp, plus 20 core RefSeq nuclear transcripts \texttt{NUC-101}--\texttt{NUC-120} = 52,925 bp drawn from 48 verified RefSeq orthologs)~\cite{Soghigian2023,Reidenbach2009,Neafsey2015,Poinar2019}, we execute two complementary empirical chains: \textbf{Pipeline A} (16,032 bp MAFFT v7.526 Total-Evidence Supermatrix + IQ-TREE v2.4.0 + LSD2 + PAML v4.10.10, Figure 1, $\ln L = -152,134.21$, $\text{BIC} = 306,476.00$) and \textbf{Pipeline B} (4,140 bp MAFFT v7.526 100-Species Mitochondrial Supermatrix + IQ-TREE v2.4.0 $-B\ 1000$, Figure 5, $\ln L = -67,577.15$, $\text{BIC} = 136,961.58$)~\cite{daSilva2020,Foster2017}.

\section{Stage 3 ModelTest Compositional Gate \& Discovery of GenBank Sample Swap \texttt{PQ585799.1}}
Across the 100-species mitochondrial 4-gene CDS supermatrix (\texttt{COX1+COB+COX2+COX3} = 4,140 bp), raw third-codon positions violate the compositional homogeneity threshold ($\text{RCFV}_{\text{Pos3}} = 0.0737 > 0.0450$). Purine/pyrimidine third-codon recoding (\texttt{RY3}) reduces RCFV by 76.3\% to $0.0175$. Dual-pass SHA-256 blinded validation ($Z_{\text{blind}} = +16.41$ SD vs.\ $Z_{\text{ctx}} = +16.10$ SD; $D_{\text{ctx}} = 0.31$ SD) and IQ-TREE v2.4.0 phylogenetics revealed that GenBank accession \texttt{PQ585799.1} (deposited as \textit{Toxorhynchites rutilus}~\cite{Ali2025}) nests inside \textit{Culex} (100\% UFBoot2) and is 100.0000\% identical across all 14,500 bp of coding/rRNA/tRNA sequence to \textit{Culex tarsalis} (\texttt{PQ585801.1}) from the same study.

\section{PAML v4.10.10 Codon Selection on \textit{Toxorhynchites rutilus} \& Vector Control Synthesis}
PAML v4.10.10 (\texttt{codeml}) 3-branch selection modeling~\cite{Yang2007} across 48 verified RefSeq nuclear orthologs (including all 6 authentic \textit{Toxorhynchites rutilus septentrionalis} RefSeq mRNAs from \texttt{GCF\_029784165.1})~\cite{Donald2020} indicates $2.11\times$ relaxed purifying selection on salivary \textit{D7L1} ($\omega_{\text{Toxo}} = 0.2138$ vs.\ $\omega_{\text{Culicinae}} = 0.1011$)~\cite{Calvo2006,Mans2007} and $2.09\times$--$4.48\times$ relaxed purifying selection on secreted salivary apyrase ($\omega_{\text{Toxo}} = 0.1025$ vs.\ $0.0229$--$0.0491$)~\cite{Champagne1995}, contrasted with extreme purifying conservation of housekeeping membrane 5$'$-nucleotidase ($\omega_{\text{Toxo}} = 0.0011$), olfactory \textit{Orco} ($\omega_{\text{Toxo}} = 0.0284$), and autogenous yolk-precursor \textit{Vitellogenin-A1} ($\omega_{\text{Toxo}} = 0.0929$ vs.\ $0.1457$ in biting Culicinae, $p < 10^{-16}$).

\begin{thebibliography}{13}
\bibitem{White2026RuNIE} B.~P. White, "RuNIE (Rust Native Inference Engine): An Open-Source Architecture for Autonomous Bioinformatics Pipeline Synthesis, Blinded Validation, and Bayesian Scope-Depth Auditing," \emph{CRV Science Labs / Zenodo}, 2026. \url{https://doi.org/10.5281/zenodo.23969275}
\bibitem{Soghigian2023} J.~Soghigian, C.~Sither, S.~A. Justi, \emph{et~al.}, "Phylogenomics reveals the evolutionary history of mosquitoes and the single origin of hematophagy in Culicidae," \emph{Current Biology}, vol.~33, no.~22, pp.~4864--4874, 2023. \url{https://doi.org/10.1016/j.cub.2023.09.060}
\bibitem{Reidenbach2009} K.~R. Reidenbach, S.~Cook, M.~A. Bertone, R.~E. Harbach, B.~M. Wiegmann, and N.~J. Besansky, "Phylogenetic analysis and temporal diversification of mosquitoes (Diptera: Culicidae) based on nuclear genes and morphology," \emph{BMC Evolutionary Biology}, vol.~9, p.~298, 2009. \url{https://doi.org/10.1186/1471-2148-9-298}
\bibitem{Neafsey2015} D.~E. Neafsey, R.~M. Waterhouse, M.~R. Abai, \emph{et~al.}, "Highly evolvable malaria vectors: The genomes of 16 \emph{Anopheles} mosquitoes," \emph{Science}, vol.~347, no.~6217, p.~1258522, 2015. \url{https://doi.org/10.1126/science.1258522}
\bibitem{daSilva2020} A.~F. da~Silva, L.~C. Machado, M.~B. de~Paula, \emph{et~al.}, "Mitogenomics and evolutionary timescale of Neotropical Anophelinae and Culicinae mosquitoes," \emph{Parasites \& Vectors}, vol.~13, p.~344, 2020. \url{https://doi.org/10.1186/s13071-020-04215-z}
\bibitem{Donald2020} C.~L. Donald, P.~Siriyasatien, and A.~Kohl, "\emph{Toxorhynchites} species: Biology, autogeny, and biological control potential against container-breeding mosquito vectors," \emph{Insects}, vol.~11, no.~11, p.~747, 2020. \url{https://doi.org/10.3390/insects11110747}
\bibitem{Calvo2006} E.~Calvo, B.~J. Mans, J.~F. Andersen, and J.~M.~C. Ribeiro, "Function and structure of a mosquito salivary protein that binds biogenic amines and eicosanoids," \emph{Journal of Biological Chemistry}, vol.~281, no.~4, pp.~1935--1942, 2006. \url{https://doi.org/10.1074/jbc.M510359200}
\bibitem{Mans2007} B.~J. Mans, E.~Calvo, J.~M.~C. Ribeiro, and J.~F. Andersen, "The crystal structure of D7r4, a salivary biogenic amine-binding protein from the malaria mosquito \emph{Anopheles gambiae}," \emph{Journal of Biological Chemistry}, vol.~282, no.~50, pp.~36626--36633, 2007. \url{https://doi.org/10.1074/jbc.M706410200}
\bibitem{Champagne1995} D.~E. Champagne, C.~T. Smartt, J.~M.~C. Ribeiro, and A.~A. James, "The salivary gland-specific apyrase of the mosquito \emph{Aedes aegypti} is a member of the 5$'$-nucleotidase family," \emph{Proceedings of the National Academy of Sciences}, vol.~92, no.~3, pp.~694--698, 1995. \url{https://doi.org/10.1073/pnas.92.3.694}
\bibitem{Poinar2019} G.~Poinar, T.~J. Zavortink, and A.~Brown, "\emph{Priscoculex burmanicus} n. gen., n. sp. (Diptera: Culicidae: Anophelinae) from mid-Cretaceous Burmese amber," \emph{Historical Biology}, vol.~32, no.~9, pp.~1157--1162, 2019. \url{https://doi.org/10.1080/08912963.2019.1570185}
\bibitem{Foster2017} P.~G. Foster, T.~M.~P. de~Oliveira, E.~S. Bergo, \emph{et~al.}, "Phylogeny of Anophelinae using mitochondrial protein coding genes," \emph{Royal Society Open Science}, vol.~4, no.~11, p.~170758, 2017. \url{https://doi.org/10.1098/rsos.170758}
\bibitem{Yang2007} Z.~Yang, "PAML 4: Phylogenetic Analysis by Maximum Likelihood," \emph{Molecular Biology and Evolution}, vol.~24, no.~8, pp.~1586--1591, 2007. \url{https://doi.org/10.1093/molbev/msm088}
\bibitem{Ali2025} R.~Ali \emph{et~al.}, "Complete mitochondrial genomes of mosquito species from Harris County, Texas (GenBank accessions PQ585799.1 and PQ585801.1)," \emph{Scientific Reports}, vol.~15, p.~22013, 2025. \url{https://doi.org/10.1038/s41598-025-06412-9}
\end{thebibliography}
\end{document}
